\chapter{The Platform}\label{ch:platform}

This chapter describes what the experiments run on: the robot and its cameras, the computers, the software layer we
built on LeRobot~0.7, and how work moves between the machines. Figure~\ref{fig:pipeline} shows the whole path of the
project, from a recorded dataset to the report; the following chapters take its stages in turn.

\begin{figure}[htbp]
  \centering
  % The whole project: record -> train -> evaluate (three ways) -> report, and what carries the work between stages.
% Shared figure colours (colour only in figures).
\definecolor{robotfill}{HTML}{EAF2FC}\definecolor{robotline}{HTML}{2A78D6}
\definecolor{serverfill}{HTML}{FDEEE7}\definecolor{serverline}{HTML}{EB6834}
\definecolor{netfill}{HTML}{F0EFEC}\definecolor{greenfill}{HTML}{E9F5EC}\definecolor{greenline}{HTML}{2F9E58}
\definecolor{purplefill}{HTML}{F1ECF9}\definecolor{purpleline}{HTML}{7A55C7}
\definecolor{ink}{HTML}{0B0B0B}\definecolor{ink2}{HTML}{52514E}\definecolor{faultc}{HTML}{C9C6BF}

\begin{tikzpicture}[font=\small, node distance=5mm and 7mm,
  st/.style={draw=#1, thick, rounded corners=2pt, fill=white, align=center, minimum width=29mm, minimum height=13mm,
             inner sep=3pt},
  arr/.style={-{Latex[length=2.2mm]}, thick, ink2}]
\node[st=robotline] (rec) {\textbf{Record}\\ leader arm, cameras\\ trim, review};
\node[st=serverline, right=of rec] (tr) {\textbf{Train}\\ 7 policies,\\ original recipes};
\node[st=greenline, right=of tr, yshift=17mm] (off) {\textbf{Offline}\\ held-out error,\\ inference cost};
\node[st=greenline, right=of tr] (rob) {\textbf{Robot}\\ trials, success,\\ the conditions};
\node[st=greenline, right=of tr, yshift=-17mm] (ben) {\textbf{Bench}\\ replay robot,\\ every fault};
\node[st=purpleline, right=of rob] (rep) {\textbf{Report}\\ tables, figures\\ from raw files};
\draw[arr] (rec) -- (tr); \draw[arr] (tr) -- (rob); \draw[arr] (tr.east) -- ++(2mm,0) |- (off.west);
\draw[arr] (tr.east) -- ++(2mm,0) |- (ben.west);
\draw[arr] (rob) -- (rep); \draw[arr] (off.east) -- ++(2mm,0) |- (rep.west); \draw[arr] (ben.east) -- ++(2mm,0) |- (rep.west);
\node[below=4mm of rec, font=\scriptsize, text=ink2, align=center] {datasets on the Hub\\ \texttt{recorded.yaml} in git};
\node[below=4mm of tr, font=\scriptsize, text=ink2, align=center] {models on the Hub\\ (best and final checkpoint)};
\node[below=2mm of ben, font=\scriptsize, text=ink2, align=center] {sessions in the\\ results store};
\end{tikzpicture}

  \caption{The path of the project. Datasets are recorded and trimmed on the laptop, policies are trained on the GPU
           servers, and every trained policy is evaluated in three ways: offline, on the robot, and on the replay
           bench. The report is rebuilt from the raw result files.}
  \label{fig:pipeline}
\end{figure}

\section{Hardware}\label{sec:hardware}

\begin{table}[htbp]
  \centering
  \small
  \begin{tabularx}{\textwidth}{l X}
    \toprule
    Component & Description \\
    \midrule
    Arm & SO-101 follower arm~\cite{so101}: six Feetech servos, joint targets in degrees; an identical leader arm,
          moved by hand, for recording \\
    Wrist camera & Zebronics Zeb Live Pro on the gripper, $800\times600$ MJPG at 30~fps \\
    Scene camera & An Android phone over USB, bridged to a virtual video device (DroidCam), $800\times600$ at 30~fps;
                   always the same phone, on a fixed mount \\
    Robot computer & Laptop: Linux, NVIDIA RTX~4060 Laptop GPU (8~GB). Runs the arm, the cameras and the client at
                     30~Hz; the robot computer of every official experiment \\
    Edge computer & Raspberry Pi~5 (8~GB, no GPU), as an alternative robot computer (Chapter~\ref{ch:edge}) \\
    GPU server 1 & NVIDIA A100 80~GB PCIe on a shared server (driver for CUDA $\le$ 12.8) \\
    GPU server 2 & NVIDIA H200 NVL 141~GB on a shared server, inside a container (driver for CUDA $\le$ 13.0) \\
    Network & Laptop to server over WiFi and the local network; the Pi reaches the laptop through its hotspot \\
    \bottomrule
  \end{tabularx}
  \caption{Hardware.}\label{tab:hardware}
\end{table}

Table~\ref{tab:hardware} lists the hardware and Figures~\ref{fig:setup} and~\ref{fig:robot} show it. The two cameras
have fixed roles: the wrist camera sees the object and the gripper at close range; the scene camera sees the whole
task area. Both are read at $800\times600$, and every request to a policy server carries both images uncompressed
($2\times800\times600\times3$~bytes $\approx$~2.9~MB) together with the six joint positions. The phone is bridged
to a virtual video device at every boot so that it appears to LeRobot like any other camera; the arms, their serial
ports and the calibration files are described in a versioned \emph{rig} file, so every recording and every trial
records exactly which hardware it used.

\begin{figure}[htbp]
  \centering
  \photo[0.9\textwidth]{setup.jpg}{The whole setup: the SO-101 arm, the task board, the cameras and the laptop.}
  \caption{The experimental setup.}
  \label{fig:setup}
\end{figure}

\begin{figure}[htbp]
  \centering
  \begin{subfigure}[t]{0.48\textwidth}
    \centering
    \photo[\textwidth]{arm_wrist_camera.jpg}{The follower arm with the wrist camera.}
    \caption{Follower arm and wrist camera.}
  \end{subfigure}\hfill
  \begin{subfigure}[t]{0.48\textwidth}
    \centering
    \photo[\textwidth]{scene_camera.jpg}{The scene camera (phone) and its view of the task.}
    \caption{Scene camera.}
  \end{subfigure}
  \caption{The robot and its cameras.}
  \label{fig:robot}
\end{figure}

\section{Software: a layer on LeRobot 0.7}\label{sec:software}

The project is built on LeRobot~0.7~\cite{lerobot}, used as a pinned dependency and not modified. Everything of our
own lives in a separate layer on top of it, so that LeRobot's training, policy server, robot drivers and safety checks
are used exactly as they ship, and every change we make is in our own code where it can be tested. The layer has
four parts:
\begin{itemize}
\item \textbf{The registry}: plain YAML files that describe everything an experiment depends on (next section).
\item \textbf{The \texttt{lab} command}: one entry point that records, trims, trains, serves, runs and evaluates from
      the registry, and builds the exact LeRobot command for each step.
\item \textbf{Two LeRobot plugins}, found by LeRobot through its plugin mechanism: GR00T~N1.5, which LeRobot removed
      after version~0.5.1, ported as a policy plugin with guided RTC (Section~\ref{sec:rtcmethod}); and a
      \emph{replay robot}, a robot driver that plays a recorded episode as its camera images and joint states, which
      the bench uses in place of the arm (Section~\ref{sec:bench-design}).
\item \textbf{The control layer}: our \drtc{} scheduler, the smoothing methods and the fault injector, which run inside
      LeRobot's remote-inference engine (Chapter~\ref{ch:control}), and the evaluation framework
      (Chapter~\ref{ch:evaluation}).
\end{itemize}
The software runs on Linux on every machine: Python~3.12, PyTorch~2.14.1 built for each server's CUDA version, and
LeRobot's optional dependencies for the robot, the cameras and the policies. The layer is covered by
\R{software/tests} automated tests, which run on the laptop without a GPU or a robot.

\section{The registry}\label{sec:registry}

Every object an experiment depends on is a small file in the repository, and every file refers to the others by name:

\begin{table}[htbp]
  \centering
  \small
  \begin{tabularx}{\textwidth}{l X}
    \toprule
    File & What it fixes \\
    \midrule
    Rig & the arms (serial ports per machine, calibration files), the cameras (device, role, resolution) \\
    Dataset & the Hub repository, pinned to one revision with its episode and frame counts, the joint units, the task
              instruction, the held-out split, and which recording it was trimmed from \\
    Recipe & how a policy is trained: its authors' settings, unchanged (Section~\ref{sec:recipes}) \\
    Model & a recipe applied to a dataset; its Hub repository and status \\
    Host & a machine: its role, GPUs, login, folders, secrets file and CPU budget for parallel trainings \\
    Condition & one way of running the robot side and the network: scheduler, smoothing, faults
                (Chapter~\ref{ch:control}) \\
    Evaluation & one study: the models, conditions, trials, success rule and time limit, frozen before the first
                 trial (Chapter~\ref{ch:evaluation}) \\
    \bottomrule
  \end{tabularx}
  \caption{The registry: what each kind of file fixes.}\label{tab:registry}
\end{table}

Because a model names its recipe and dataset, and the dataset names its revision, a model's name determines exactly
what it was trained on and how. A condition is combined with the shared defaults, the rig and the model, so the same condition
file is used for every policy. Generated files are made by small scripts in the registry itself; the conditions of the
main experiments (Chapter~\ref{ch:design}) are generated this way, so their systems and faults are identical
everywhere.

\section{Machines and how work moves between them}\label{sec:machines}

\begin{figure}[htbp]
  \centering
  % Machines and the three places that carry work between them.
% Shared figure colours (colour only in figures).
\definecolor{robotfill}{HTML}{EAF2FC}\definecolor{robotline}{HTML}{2A78D6}
\definecolor{serverfill}{HTML}{FDEEE7}\definecolor{serverline}{HTML}{EB6834}
\definecolor{netfill}{HTML}{F0EFEC}\definecolor{greenfill}{HTML}{E9F5EC}\definecolor{greenline}{HTML}{2F9E58}
\definecolor{purplefill}{HTML}{F1ECF9}\definecolor{purpleline}{HTML}{7A55C7}
\definecolor{ink}{HTML}{0B0B0B}\definecolor{ink2}{HTML}{52514E}\definecolor{faultc}{HTML}{C9C6BF}

\begin{tikzpicture}[font=\small, node distance=6mm and 9mm,
  m/.style={draw=#1, thick, rounded corners=2pt, fill=white, align=center, minimum width=30mm, minimum height=13mm},
  s/.style={draw=ink2, thick, rounded corners=2pt, fill=netfill, align=center, minimum width=34mm, minimum height=11mm},
  arr/.style={{Latex[length=2mm]}-{Latex[length=2mm]}, thick, ink2}]
\node[m=robotline] (lap) {\textbf{Laptop}\\ recording, robot computer\\ RTX 4060 8 GB};
\node[m=robotline, below=of lap] (pi) {\textbf{Raspberry Pi 5}\\ edge robot computer};
\node[s, right=18mm of lap, yshift=10mm] (git) {\textbf{git} (this repository)\\ registry, eval plans};
\node[s, below=of git] (hub) {\textbf{Hugging Face Hub}\\ datasets, models};
\node[s, below=of hub] (res) {\textbf{Results store}\\ (private Hub dataset)};
\node[m=serverline, right=18mm of git, yshift=-6mm] (a100) {\textbf{A100} 80 GB\\ train, serve, bench};
\node[m=serverline, below=of a100] (h200) {\textbf{H200} 141 GB\\ train, serve, bench\\ (container)};
\foreach \a in {lap,pi} { \draw[arr] (\a.east) -- (git.west); \draw[arr] (\a.east) -- (hub.west); \draw[arr] (\a.east) -- (res.west); }
\foreach \b in {a100,h200} { \draw[arr] (\b.west) -- (git.east); \draw[arr] (\b.west) -- (hub.east); \draw[arr] (\b.west) -- (res.east); }
\draw[-{Latex[length=2mm]}, thick, densely dashed, robotline] (lap.north) to[bend left=25] node[above, font=\scriptsize]{observations / chunks (SSH tunnel)} (a100.north);
\end{tikzpicture}

  \caption{The machines and the three places that carry work between them. Code and plans travel through git,
           datasets and models through the Hugging Face Hub, and every official result through a private results
           store. During a robot session the laptop reaches the model server through an SSH tunnel.}
  \label{fig:machines}
\end{figure}

Work is never copied between machines by hand (Figure~\ref{fig:machines}). Git carries the code and the registry, so a
GPU server learns about a new dataset or an evaluation plan with a \texttt{git pull}. The Hugging Face Hub carries the
datasets and every checkpoint of every trained model. A private dataset repository on the Hub serves as the
\emph{results store}: every official session is pushed there when it ends, and every report is rebuilt from it.

Both GPU servers are shared with other users and cannot be changed, so the software adapts to them: each gets the
PyTorch build its driver supports, long jobs run in terminal multiplexer sessions that survive a logout, the H200's
work runs inside a container started the way the server's guide prescribes, and nothing outside the project's folder
is ever written or deleted. During a robot session the laptop reaches the model server through one SSH connection
that forwards a port for every served model; on the H200 the connection reaches into the container.
